\documentclass{handout}
\usepackage{handoutSetup}\usepackage{handoutShortcuts}  

\begin{document}
\handoutHeader

\begin{verbatimwrite}{\jobname.title}
Social Security and Capital Accumulation
\end{verbatimwrite}

\handoutNameMake


Consider a household with a 2-period lifetime, whose optimization 
problem is\opt{MarginNotes}{\marginpar{\tiny Written in the most general way possible.  Last line indicates assumption that no labor income is earned in period 2.}}
\begin{equation}\begin{gathered}\begin{aligned}
        \vFunc(\bRat_{1,t}) & =  \max_{\{\cRat_{1,t}\}} ~~ \uFunc(\cRat_{1,t}) + \Discount \uFunc(\cRat_{2,t+1})  \\
         & \mbox{s.t.}    \nonumber 
\\      \yRat_{1,t} & =  (\Wage_{1,t}-\taxNetTrans_{1,t})
\\      \aRat_{1,t} & =  \yRat_{1,t} - \cRat_{1,t}
\\  \yRat_{2,t+1} & =  (\Wage_{2,t+1}-\taxNetTrans_{2,t+1}) = -\taxNetTrans_{2,t+1}
\\  \cRat_{2,t+1} & =  \Rfree_{t+1} \aRat_{1,t}+\yRat_{2,t+1}
.
\end{aligned}\end{gathered}\end{equation}

Under logarithmic utility, handout \handoutC{2PeriodLCModel} shows that the solution to this problem is
\begin{equation}\begin{gathered}\begin{aligned}
        \cRat_{1,t} & =  \left(\frac{\yRat_{1,t}+\yRat_{2,t+1}/\Rfree_{t+1}}{1+\Discount}\right).  \label{eq:c1} 
\end{aligned}\end{gathered}\end{equation}

The only role of government in this economy is to run a Social
Security program.  Suppose that initially this economy
had no Social Security system and we are interested in the effects of
introducing a Pay-As-You-Go Social Security system that is expected to
remain a constant size from generation to generation from now on:
$\taxNetTrans_{2,t+1} = -\taxNetTrans_{1,t+1}$ while $\taxNetTrans_{1,t+1} =
\taxNetTrans_{1,t}$, so that taxes are greater than transfers when young
and transfers are greater than taxes when old.

The effects of Social Security on first period consumption can be 
seen by writing out explicitly the value for $\cRat_{1,t}$ from equation \eqref{eq:c1}, \opt{MarginNotes}{\marginpar{\tiny Substitute definitions of $\yRat_{1,t}$ and $\yRat_{2,t+1}$.}}
\begin{equation}\begin{gathered}\begin{aligned}
        \cRat_{1,t} & =   \left(\Wage_{1,t}-\taxNetTrans_{1,t}-\taxNetTrans_{2,t+1}/\Rfree_{t+1}\right)/(1+\Discount)
\\       & =   \left(\Wage_{1,t}-\underbrace{\rfree_{t+1}\taxNetTrans_{1,t}/\Rfree_{t+1}}_{\mbox{\footnotesize consume less b/c poorer}}\right)/(1+\Discount)
\end{aligned}\end{gathered}\end{equation}
where the expression with the underbrace comes from the effect of
introducing a constant-sized PAYG Social Security system in handout
\handoutC{GenAcctsAndGov}.  If taxes paid when young
$\taxNetTrans_{1,t}$ are positive (as they are after the introduction
of the Social Security system) and the interest rate is positive, the expression with the underbrace is
a positive number, and since it is being subtracted from $\Wage_{1,t}$
it is clear that consumption in the first period of life will {\it
  decline} with the introduction of the Social Security
system.\opt{MarginNotes}{\marginpar{\tiny Consumption declines because
    you are poorer in a lifetime sense.}}

Does the decline in consumption mean the saving rate rises?  No - 
because saving is {\it after-tax} income minus consumption, and net 
taxes on the young have risen.  For saving we have
\begin{equation}\begin{gathered}\begin{aligned}
        \aRat_{1,t} & =  \overbrace{(\Wage_{1,t}-\taxNetTrans_{1,t})}^{\mbox{\footnotesize Less after-tax income}} - \overbrace{\cRat_{1,t}}^{\mbox{\footnotesize Lower consumption b/c poorer}}  \\
         & =  (\Wage_{1,t}(1-1/(1+\Discount)) - \taxNetTrans_{1,t}+  [r_{t+1}\taxNetTrans_{1,t}/\Rfree_{t+1}]/(1+\Discount)  \\
         & =  \Wage_{1,t}\left(\frac{\Discount}{1+\Discount}\right)-\taxNetTrans_{1,t}\left(1-\frac{r_{t+1}}{\Rfree_{t+1}(1+\Discount)}\right)  \\
         & =  \Wage_{1,t}\left(\frac{\Discount}{1+\Discount}\right)-\taxNetTrans_{1,t}\left(\frac{\Rfree_{t+1}(1+\Discount)-r_{t+1}}{\Rfree_{t+1}(1+\Discount)}\right)  \\
         & =  \Wage_{1,t}\left(\frac{\Discount}{1+\Discount}\right)-\taxNetTrans_{1,t}\left(\frac{1+\Rfree_{t+1}\Discount}{\Rfree_{t+1}(1+\Discount)}\right).
\end{aligned}\end{gathered}\end{equation}
So if $\taxNetTrans_{1,t}>0$ then saving is less than before the introduction of 
Social Security.

Now consider the implications in a \cite{diamond:olg} OLG model where saving is the source 
of capital accumulation.  Suppose there is no population growth so that
\begin{equation}\begin{gathered}\begin{aligned}
        K_{t+1} & =  \aRat_{1,t}
\\       & =  \Wage_{1,t}\left(\frac{\Discount}{1+\Discount}\right)-\taxNetTrans_{1,t}\left(\frac{1+\Rfree_{t+1}\Discount}{\Rfree_{t+1}(1+\Discount)}\right)
\\   & =  (1-\varepsilon)K^{\varepsilon}_{t}      \left(\frac{\Discount}{1+\Discount}\right)-\taxNetTrans_{1,t}\left(\frac{1+\Rfree_{t+1}\Discount}{\Rfree_{t+1}(1+\Discount)}\right)
\\   & =  \mathcal{Q} K^{\varepsilon}_{t}     -\taxNetTrans_{1,t}\left(\frac{1+\Rfree_{t+1}\Discount}{\Rfree_{t+1}(1+\Discount)}\right)
\end{aligned}\end{gathered}\end{equation}
where $\mathcal{Q}=(1-\varepsilon)\Discount/(1+\Discount)$ as before in the \handoutC{OLGModel} handout.\opt{MarginNotes}{\marginpar{\tiny Admit the slight bogusness of the figure caused by the fact that $\Rfree_{t+1}$ is a function of $k_{t+1}$ while the figure assumes that the $K_{t+1}$ curve is shifted down a constant amount.}}
Thus the capital accumulation curve is shifted down.  The dynamics of 
the introduction of Social Security are captured in the figure, under 
the assumption that the economy was at its steady-state equilibrium 
level $\bar{k}$ before the Social Security system was introduced.  The 
effect of introduction is an immediate increase in consumption, as the 
old generation spends everything it gets and the young generation 
doesn't need to do as much retirement saving as before.  Over time the 
economy will converge to its new, lower level of capital $\bar{\bar{k}}$.

% Figure is generated in the Mathematica notebook OLGModel.nb

\begin{figure}
\includegraphics[width=6in]{../Figures/SocSecAndKAccum}
        \caption{Convergence of OLG Economy After Intro of Social Security}
        \label{fig:SocSecAndKAccum}
\end{figure}


\input handoutBibMake

\end{document}
